\section{The residual pair: a physical-time core
trace}

\subsection{1. Scope and result}

This note concerns only the two linkage classes

\[
 {\cal L}_0=\{B,2A,B+C\},\qquad
 {\cal L}_1=\{0,A,C\}.                                    \tag{1.1}
\]

Each labelled reaction graph is assumed strongly connected, parallel
channels in the same direction are aggregated, and every displayed
channel rate is positive. A missing directed edge has aggregate rate
zero. All propensities are stochastic mass-action propensities.

The purpose of the argument is to avoid both stationary-only averaging
and recurrence of the raw jump chain. The chain is sampled after a
physical window of length \(T/\sqrt{q+1}\), where

\[
 q=A+2B.                                                     \tag{1.2}
\]

The fast linkage preserves \(q\). A rate-dependent proper linear
workload returns the process to a core on which \(A=O(\sqrt q)\) and
\(C=O(1)\). On the short physical window the exact fast dynamics have a
deterministic Riccati limit, while all positive \(q\)-hazards vanish.
This gives a strict, shell-uniform negative drift for the core trace.
Endpoint moments and the physical duration of the return appendage are
included below.

\begin{quote}
\textbf{Theorem 1.1.} For every choice of strongly connected
orientations and positive rates in (1.1), the minimal stochastic
mass-action CTMC is nonexplosive. Every state in every closed
irreducible population class is positive recurrent.
\end{quote}

No finite inactive-coordinate box is assumed. The finite \(C\)-bound
below is a state-dependent core produced by a proper Foster workload;
excursions outside it are retained and charged.

\subsection{2. Aggregated rates and two exact
identities}

Write the six possible rates in \({\cal L}_0\) as

\[
\begin{array}{c|cccccc}
\text{edge}&B\to2A&B\to B+C&2A\to B&2A\to B+C&B+C\to B&B+C\to2A\\
\hline
\text{rate}&x&y&s&r&t&v .
\end{array}                                                  \tag{2.1}
\]

Put \(d=t+v\). Strong connectivity gives \(d>0\). Define

\[
 \zeta={2v\over d},\qquad
 \alpha=x+{vy\over d},\qquad
 \beta=2\left(s+{rt\over d}\right).                        \tag{2.2}
\]

Both \(\alpha\) and \(\beta\) are strictly positive. For example,
\(\alpha=0\) would imply that neither \(B\) nor \(B+C\) can reach
\(2A\); the proof for \(\beta\) is the reverse argument.

Set

\[
 Z=A+\zeta C.
\]

Direct substitution gives the exact identity

\[
 {\cal L}_0Z
 =2\alpha B-\beta(A)_2
 =\alpha(q-A)-\beta A(A-1)=:P_q(A).                         \tag{2.3}
\]

The coefficient of \(BC\) cancels because \(-\zeta t+(2-\zeta)v=0\). Let
\(a_q\) be the positive root of \(P_q(a)=0\). Then

\[
 {a_q\over\sqrt{q+1}}\longrightarrow
 a_*:=\sqrt{\alpha/\beta}.                                 \tag{2.4}
\]

For \({\cal L}_1\), write its rates as \(k_{ij}\), with
\(i,j\in\{0,A,C\}\). Then

\[
 {\cal L}_1q=k_{0A}+k_{CA}C-wA,\qquad
 w:=k_{A0}+k_{AC}>0.                                       \tag{2.5}
\]

The strict positivity of \(w\) follows from strong connectivity.

\subsection{3. A proper workload and an exact
core}

The return workload is constructed algebraically. Use the conventions
\(d/v=+\infty\) if \(v=0\), and \((k_{C0}+k_{CA})/k_{CA}=+\infty\) if
\(k_{CA}=0\). Strong connectivity implies that the open intervals

\[
 I_0=\left({r\over s+r},{d\over v}\right),\qquad
 I_1=\left({k_{AC}\over k_{A0}+k_{AC}},
            {k_{C0}+k_{CA}\over k_{CA}}\right)             \tag{3.1}
\]

are nonempty and contain values arbitrarily close to one. If an endpoint
equals one, the opposite endpoint is strict; equality at both ends would
disconnect one of the three-vertex graphs.

Choose \(\lambda\in I_0\) and \(\rho\in I_1\), sufficiently close to one
that \(2\rho-\lambda>0\), and put

\[
 U=\rho A+(2\rho-\lambda)B+C.                              \tag{3.2}
\]

Define the positive constants

\[
\begin{aligned}
 c_B&=x\lambda+y,\\
 c_2&=(s+r)\lambda-r,\\
 c_{BC}&=d-v\lambda,\\
 d_A&=(k_{A0}+k_{AC})\rho-k_{AC},\\
 d_C&=(k_{C0}+k_{CA})-k_{CA}\rho,\\
 K_0&=k_{0A}\rho+k_{0C}.
\end{aligned}                                               \tag{3.3}
\]

The generator identity is

\[
 {\cal L}U
 =K_0+c_BB-c_2(A)_2-c_{BC}BC-d_AA-d_CC.                    \tag{3.4}
\]

This identity is valid at the boundary because falling factorials,
rather than ordinary monomials, occur in (3.4).

Put

\[
 D(A,B,C)=c_2(A)_2+c_{BC}BC+d_AA+d_CC                       \tag{3.5}
\]

and define

\[
 {\cal R}=\{D\le 2(K_0+c_BB+1)\}.                         \tag{3.6}
\]

Outside \({\cal R}\),

\[
 {\cal L}U\le-1-\tfrac12D.                               \tag{3.7}
\]

There are explicit constants

\[
 K_A=2+\sqrt{{2(K_0+c_B+1)\over c_2}},qquad
 K_C={2(K_0+c_B+1)\over c_{BC}}                            \tag{3.8}
\]

and a finite integer \(q_0\) such that every state of \({\cal R}\) with
\(q>q_0\) satisfies

\[
 A\le K_A\sqrt{q+1},\qquad C\le K_C.                       \tag{3.9}
\]

Indeed, (3.9) follows separately from the \((A)_2\) and \(BC\) terms in
(3.5). When \(B=0\), (3.6) bounds both \(A\) and \(C\), so all
sufficiently large states in \({\cal R}\) have \(B\ge1\).

Choose \(K\) and \(C_*\), after \(q_0\) has been fixed, so that

\[
\begin{split}
 K&>K_A+a_*+2,\qquad C_*>K_C+2,\\
 K&>1+\max_{x\in{\cal R},\ q(x)\le q_0}
          {A(x)\over\sqrt{q(x)+1}},\\
 C_*&>1+\max_{x\in{\cal R},\ q(x)\le q_0}C(x).
\end{split}                                                 \tag{3.10}
\]

and define the moving core

\[
 {\cal K}=\{(a,b,c):a\le K\sqrt{a+2b+1},\ c\le C_*\}.       \tag{3.11}
\]

The finite maxima in (3.10) exist because (3.6) bounds \(C\) when \(q\)
is bounded. Together with (3.9), they give the global inclusion

\[
 {\cal R}\subset{\cal K}.                                  \tag{3.12}
\]

Thus (3.7) is a genuine proper-workload return mechanism to
\({\cal K}\), not a truncation assumption.

The positive part of (3.4) is at most \(K(1+U)\). Localization and
Gronwall therefore give

\[
 \mathbb E_xU(X_{t\wedge\tau_n})
 \le (U(x)+K t)e^{Kt},                                     \tag{3.13}
\]

where \(\tau_n\) exits the \(U\)-sublevel \(\{U\le n\}\). Since those
sublevels are finite, (3.13) proves nonexplosion and endpoint
integrability. The same calculation applied to \(U^p\), using the strict
negative coefficients at the two quadratic sources, gives finite
polynomial moments on every bounded physical interval.

\subsection{4. The short physical
window}

For a starting state in \({\cal K}\), write \(N=q(x)\),
\(S_N=\sqrt{N+1}\), and fix \(T>0\). During the physical interval

\[
 h_N={T\over S_N},                                         \tag{4.1}
\]

use accelerated time \(\tau=S_Nt\) and set

\[
 z_N(\tau)={A(\tau/S_N)+\zeta C(\tau/S_N)\over S_N}.        \tag{4.2}
\]

Let \(\Phi_\tau(u)\) denote the solution at time \(\tau\) of

\[
 \dot z=\alpha-\beta z^2,\qquad z(0)=u.                    \tag{4.3}
\]

The following is the quantitative finite-time statement needed below.

\begin{quote}
\textbf{Lemma 4.1 (uniform fast window).} Uniformly over starting states
in \({\cal K}\) with \(q=N\), as \(N\to\infty\), \[
\mathbb E_x\sup_{0\le\tau\le T}
  |z_N(\tau)-\Phi_\tau(z_N(0))|^2\longrightarrow0.        \tag{4.4}
\] For every integer \(p\ge1\), \[
\sup_N\sup_{0\le\tau\le T}
  \mathbb E_x[C(\tau/S_N)^p]<\infty,                     \tag{4.5}
\] and the number \(J_N\) of \(q\)-changing reactions in the window has
uniformly bounded \(p\)-th moment. Finally, \[
\sup_{x\in{\cal K}:\,q(x)=N}\left|
\mathbb E_x[q(X_{h_N})-N]
+w\int_0^T\Phi_\tau(z_N(0))\,d\tau\right|
\longrightarrow0.                                      \tag{4.6}
\]
\end{quote}

\subsubsection{Proof of Lemma 4.1}

Let \(\vartheta_N\) be the first accelerated time at which
\(q\notin[N/2,2N]\) or \(A>2KS_N\). On \([0,T\wedge\vartheta_N]\),
\(B\ge N/8\) for all large \(N\). The physical \(C\)-birth rate is at
most

\[
 yB+r(A)_2+k_{0C}+k_{AC}A\le b_CN,                         \tag{4.7}
\]

where one may take \[
 b_C=y+8rK^2+k_{0C}+2Kk_{AC}+1.
\] The physical \(C\)-death rate contains

\[
 dBC\ge\delta_CN C,\qquad \delta_C=d/8.                    \tag{4.8}
\]

Consequently, up to \(\vartheta_N\), \(C\) is stochastically dominated
by the immigration--death chain \(\bar C\) with immigration rate
\(b_CN\) and per-particle death rate \(\delta_CN\), started at
\(C(0)\le C_*\). At every physical time \(t\),

\[
\bar C(t)\ \stackrel d=\ {\rm Bin}\!\left(C(0),e^{-\delta_CNt}\right)
 +{\rm Pois}\!\left({b_C\over\delta_C}
             (1-e^{-\delta_CNt})\right).                  \tag{4.9}
\]

In particular, for every \(\eta>0\),

\[
 \sup_N\sup_{\tau\le T}
 \mathbb E_x\exp\{\eta C((\tau\wedge\vartheta_N)/S_N)\}
 \le
 \exp\!\left\{\eta C_*+{b_C\over\delta_C}(e^\eta-1)\right\}
 =:M_\eta.                                                \tag{4.10}
\]

This is a transient coupling from every core point, not a
stationary-start estimate.

By (2.3), the predictable drift of (4.2) is

\[
 {P_q(A)\over N}+O(N^{-1/2})(1+C)
 =\alpha-\beta z_N^2+o_{L^p}(1)                            \tag{4.11}
\]

uniformly up to \(\vartheta_N\), for every fixed \(p\). Here we used
\(A/S_N=z_N-\zeta C/S_N\), (4.10), and the fact that the \(q\)-changing
intensity on accelerated time is bounded by \(K_1(1+C)\). Its counting
process therefore has moments of every fixed order bounded uniformly in
\(N\).

If \(M_N\) denotes the martingale part of \(z_N\), its jumps have size
at most \(2/S_N\), and its predictable quadratic variation is bounded by

\[
 \langle M_N\rangle_T
 \le {K_2\over S_N^3}\int_0^{T\wedge\vartheta_N}
       \{N+(A)_2+NC\}\,d\tau.                             \tag{4.12}
\]

The Burkholder--Davis--Gundy inequality, including its jump term, and
(4.10) imply, for every integer \(m\ge1\),

\[
 \mathbb E_x\sup_{\tau\le T\wedge\vartheta_N}|M_N(\tau)|^{2m}
 \le C_mS_N^{-m}.                                        \tag{4.13}
\]

The same estimate applied to the integral equation (4.11), followed by
Gronwall on the compact stopped range, gives, for every \(\epsilon>0\),

\[
 \mathbb P_x\!\left\{
 \sup_{\tau\le T\wedge\vartheta_N}
 |z_N(\tau)-\Phi_\tau(z_N(0))|>\epsilon\right\}
 \le C_{\epsilon,m}S_N^{-m}.                              \tag{4.14}
\]

The Riccati solutions starting in \([0,K+1]\) stay a fixed distance
below \(2K\). Exiting the \(q\)-band requires at least \(N/2\) bounded
\(q\)-jumps. The all-order count bounds and (4.14) therefore yield

\[
 \mathbb P_x\{\vartheta_N\le T\}\le C_mS_N^{-m}
 \quad\hbox{for every }m.                                 \tag{4.15}
\]

On this exceptional event, the polynomial endpoint estimates from (3.13)
and Hölder's inequality let us remove the stopping in (4.10)--(4.14).
This proves (4.4), (4.5), and the moment assertion for \(J_N\).

On accelerated time, (2.5) now gives

\[
 \mathbb E[q(X_{h_N})-N]
 =\int_0^T\mathbb E\left[
 {k_{0A}+k_{CA}C\over S_N}-w{A\over S_N}\right]d\tau.      \tag{4.16}
\]

The positive terms vanish by (4.5), and \(A/S_N=z_N-\zeta C/S_N\).
Equation (4.4) gives (4.6). \(\square\)

The Riccati flow is order preserving. Its smallest solution over all
nonnegative starting values is the solution from zero,

\[
 z_0(\tau)=a_*\tanh(\sqrt{\alpha\beta}\,\tau).
\]

Consequently

\[
 I_T:=\inf_{z(0)\ge0}\int_0^Tz(\tau)d\tau
 ={1\over\beta}\log\cosh(\sqrt{\alpha\beta}\,T)>0.        \tag{4.17}
\]

Lemma 4.1 therefore gives, uniformly on the core,

\[
 \mathbb E_x[q(X_{h_N})-N]\le-\tfrac34wI_T                \tag{4.18}
\]

for all sufficiently large \(N\).

\subsection{5. Returning to the core costs (o(1))
workload}

Let \(Y=X_{h_N}\), and from \(Y\) continue the original, unmodified CTMC
until its first entrance into \({\cal K}\); call the additional time
\(\sigma_N\). If \(Y\in{\cal K}\), set \(\sigma_N=0\).

\begin{quote}
\textbf{Lemma 5.1 (return appendage).} Uniformly for core starting
states, \[
\mathbb E\sigma_N=o(1),\qquad
\mathbb E[(q(X_{h_N+\sigma_N})-q(Y))^+]=o(1).             \tag{5.1}
\] All endpoint variables in (5.1) are integrable, and the full
increment \(q(X_{h_N+\sigma_N})-q(Y)\) has a uniformly bounded second
moment.
\end{quote}

\subsubsection{Proof of Lemma 5.1}

Because \(K>a_*\), order preservation of the Riccati flow gives

\[
 \Delta:=K-\sup_{0\le u\le K}\Phi_T(u)
        =K-\Phi_T(K)>0.                                   \tag{5.2}
\]

Put \(\delta=\Delta/4\). Equations (4.10), (4.14), and (4.15), together
with \(q(Y)-N=O_{L^p}(1)\), imply that for every integer \(m\ge1\),

\[
\begin{split}
 \mathbb P_x(G_N^c)&\le C_mS_N^{-m},\\
 G_N:=\bigl\{&
 3N/4\le q(Y)\le5N/4,\quad
 A(Y)\le(K-2\delta)\sqrt{q(Y)+1},\quad
 C(Y)\le S_N^{1/4}\bigr\}.
\end{split}                                                \tag{5.3}
\]

On the no-stop part of \(G_N\), \(C(Y)\) has the uniform exponential
moment (4.10). The exceptional estimate in (5.3) is super-polynomial;
the stronger \(\exp\{-cS_N\}\) estimate is not needed.

Suppose now that \(Y=y\in G_N\) and \(C(y)>C_*\). Stop the cleanup at
the first of

\[
 \gamma=\inf\{t:C(t)\le C_*\},\qquad
 \chi=\inf\{t:A(t)>K\sqrt{q(t)+1}\ {\rm or}\
                       q(t)\notin[N/2,2N]\}.              \tag{5.4}
\]

Before \(\gamma\wedge\chi\), the bounds (4.7)--(4.8) remain valid. After
the physical-time change \(u=Nt\), the dominating chain in (4.9) has
constant immigration \(b_C\) and per-particle death \(\delta_C\). The
explicit binomial--Poisson representation, followed by fixed-length
geometric trials after the initial population has contracted, gives, for
every integer \(m\ge1\),

\[
\begin{split}
 \mathbb E_y(\gamma\wedge\chi)^m
 &\le {C_m[1+\log(1+C(y))]^m\over N^m},\\
 \mathbb E_y\left(\int_0^{\gamma\wedge\chi}C(t)\,dt\right)^m
 &\le {C_m[1+C(y)]^m\over N^m}.
\end{split}\tag{5.5}
\]

For completeness, these same bounds control the probability of the stop
\(\chi\). Until \(\gamma\wedge\chi\), reactions that can change \(A\),
other than reactions consuming one \(C\), have total rate at most
\(K_3N\). The number of reactions consuming \(C\) is pathwise at most
\(C(y)\) plus the number of reactions producing \(C\). The latter has
intensity at most \(b_CN\). If \(V_{A,q}\) denotes the total variation
of \((A,q)\) before \(\gamma\wedge\chi\), thinning the corresponding
Poisson clocks and using (5.5) gives

\[
 \mathbb E_y V_{A,q}^m
 \le C_m\,\mathbb E_y
       \{1+C(y)+N(\gamma\wedge\chi)\}^m
 \le C_m\{1+C(y)+\log(1+C(y))\}^m.                        \tag{5.6}
\]

On \(G_N\), \(C(y)\le S_N^{1/4}\), whereas the \(A\)-boundary is at
least \(\delta S_N\) away and leaving the \(q\)-band requires order
\(N\) bounded jumps. Markov's inequality applied to an arbitrarily high
moment in (5.6) therefore gives

\[
 \sup_{y\in G_N}\mathbb P_y\{\chi<\gamma\}
 \le C_mS_N^{-m}\quad\hbox{for every }m.                  \tag{5.7}
\]

On \(\{\gamma<\chi\}\) the endpoint lies in \({\cal K}\). The positive
\(q\)-jump intensity is \(k_{0A}+k_{CA}C\), so (5.5) and the exponential
moment of \(C(Y)\) give

\[
 \mathbb E[\gamma;\gamma<\chi]=O(N^{-1}),\qquad
 \mathbb E N_q^+(\gamma;\gamma<\chi)=O(N^{-1}),            \tag{5.8}
\]

where \(N_q^+\) counts positive \(q\)-jumps. The negative intensity is
\(wA\le 2wKS_N\) before \(\chi\); (5.5) therefore also gives uniformly
bounded second moments for the full \(q\)-increment on the typical
cleanup.

It remains to charge \(G_N^c\) and \(\{\chi<\gamma\}\). By the global
inclusion (3.12), (3.7) holds everywhere outside \({\cal K}\). If
\(\tau_{\cal K}\) is the core hitting time, localized Dynkin estimates
give

\[
 \mathbb E_y\tau_{\cal K}
 +{1\over2}\mathbb E_y\int_0^{\tau_{\cal K}}D(X_t)\,dt
 \le U(y).                                                 \tag{5.9}
\]

Both positive and negative \(q\)-intensities are bounded by
\(K_4(1+D)\). Moreover, outside \({\cal K}\) there is a \(\kappa>0\)
such that \(D\ge\kappa(1+U)\): if \(C>C_*\), use the \(BC,A,C\) terms in
\(D\); if \(C\le C_*\), then \(A>K\sqrt{q+1}\), so the \((A)_2\) term
bounds \(q\) and hence \(U\). Also the total reaction intensity is at
most \(K_5(1+B+D)\), and outside \({\cal R}\) the \(B\)-term is bounded
by \(K_6D\). All jumps of \(U\) are bounded by a rate-dependent
constant. Hence, for each integer \(p\ge2\), the bounded-jump Taylor
formula gives, outside \({\cal K}\),

\[
\begin{split}
 {\cal L}(1+U)^p
 &=p(1+U)^{p-1}{\cal L}U+R_p,\\
 |R_p|&\le C_pD(1+U)^{p-2},\\
 {\cal L}(1+U)^p
 &\le C_p-c_pD(1+U)^{p-1}.
\end{split}\tag{5.10}
\]

For the last line, absorb the remainder when \(U\) exceeds a
\(p\)-dependent constant and enlarge \(C_p\) on the remaining finite
sublevel. Localized Dynkin's formula applied successively to
\((1+U)^j\), \(1\le j\le2p\), bounds the moments of \(\tau_{\cal K}\)
and of \(\int_0^{\tau_{\cal K}}(1+D)\,dt\). The counting-martingale
inequality then gives

\[
 \mathbb E_y\!\left[
 \tau_{\cal K}^p+
 |q(X_{\tau_{\cal K}})-q(y)|^{2p}\right]
 \le C_p(1+U(y))^{r_p}                                    \tag{5.11}
\]

for a finite integer \(r_p\) depending only on \(p\) and the rates.

Over the preceding window, (3.13) gives polynomial moments
\(\mathbb E_x(1+U(Y))^p\le C_p(1+N)^{C_p}\). At the cleanup stop the
same bound follows from (5.5)--(5.6). Hölder's inequality, (5.3), (5.7),
and arbitrarily high choices of \(m\) make the duration, positive
\(q\)-cost, and second-moment contribution of both exceptional sets
\(o(1)\). Together with (5.8), this proves (5.1). \(\square\)

Combining (4.18) and Lemma 5.1 gives a constant

\[
 \varepsilon={1\over2}wI_T>0                               \tag{5.12}
\]

and an integer \(N_0\) such that the complete core-to-core episode obeys

\[
 \mathbb E_x[q(X_\tau)-q(x)]\le-\varepsilon,qquad
 x\in{\cal K},\ q(x)>N_0.                                 \tag{5.13}
\]

Its expected physical duration is uniformly finite (in fact
\(T/S_N+o(1)\)), and its endpoint \(q\)-increment has uniformly bounded
second moment.

\subsection{6. Trace Foster argument and positive
recurrence}

Start outside \({\cal K}\) by running until the core is hit. Since
\({\cal R}\subset{\cal K}\), (3.7), localization, and Dynkin's formula
show that this time has finite expectation. At a core state with
\(q>N_0\), run the complete episode of Sections 4--5. Stop the trace
once \(q\le N_0\).

Let \(Q_n\) be the shell values of the resulting core trace and
\(\nu=\inf\{n:Q_n\le N_0\}\). Summing (5.13) only through
\(m\wedge\nu\), for finite \(m\), is legitimate by the endpoint moment
bounds. Since \(q\ge0\),

\[
 \varepsilon\,\mathbb E_x(m\wedge\nu)\le q(x),
 \qquad
 \mathbb E_x\nu\le {q(x)\over\varepsilon}.                 \tag{6.1}
\]

The episode durations are uniformly integrable, so the physical time to
hit

\[
 F={\cal K}\cap\{q\le N_0\}                                \tag{6.2}
\]

has finite expectation. The set \(F\) is finite because \(q\) bounds
\(A,B\) and the core bounds \(C\).

Fix a closed irreducible population class \(\Gamma\). The preceding
hitting argument is pathwise inside \(\Gamma\), so
\(F_\Gamma=F\cap\Gamma\) is nonempty. Record successive visits to
\(F_\Gamma\). The finite trace chain is irreducible on its visited
subset. From a nonabsorbing point of \(F_\Gamma\), take one ordinary
jump and apply the preceding finite-mean hitting bound to its finitely
many successors; an absorbing singleton is already positive recurrent.
Thus the finite trace has finite mean positive return time, and at least
one population state in \(\Gamma\) has finite mean positive physical
return time. Irreducibility gives positive recurrence of the entire
closed class. Nonexplosion was proved in (3.13), completing the proof of
Theorem 1.1.

\subsection{7. Adversarial checks}

\begin{enumerate}
\def\labelenumi{\arabic{enumi}.}
\item
  \textbf{No stationary-start assumption.} The Poisson shell law
  motivates the scale but is not used in Lemma 4.1. The \(C\)-moment
  estimate (4.9) and the Riccati limit follow from stopped generator
  identities from every core point.
\item
  \textbf{Huge initial \(C\).} Such a state is outside the core. It is
  returned by the proper workload (3.7). A large \(C\) produced at the
  end of a fast window is charged by (5.5)--(5.11); it is not discarded
  as a tail-box error.
\item
  \textbf{Boundary \(B=0\).} The bad set (3.6) is finite on \(B=0\),
  because the negative \((A)_2\) and \(C\) terms remain. Hence no large
  shell with \(B=0\) is silently included in the Riccati core.
\item
  \textbf{Positive \(C\to A\) events.} Their rate appears explicitly in
  (2.5). It vanishes on the accelerated core window by (4.5), and its
  full occupation during the return appendage is charged in (5.5),
  (5.8), and (5.9).
\item
  \textbf{Fast neutral jump counts.} No estimate counts \({\cal L}_0\)
  jumps. The proof uses physical quadratic variation over a prescribed
  physical interval and a physical hitting-time appendage.
\item
  \textbf{Endpoint integrability.} It follows from the proper linear
  workload and localized moment estimates, before any trace telescoping
  is used.
\item
  \textbf{Rate and orientation degeneracies.} Missing edges are allowed.
  The only strict constants used are
  \(\alpha,\beta,d,c_2,c_{BC},d_A,d_C,w\), all of which are forced
  positive by strong connectivity and the explicit interval choices
  (3.1).
\end{enumerate}
